\scanpage{117}\mspage{59}

Ceci vu, on peut se poser la question si $A_\lambda$, déduit de $A_0$ via $\lambda$ comme dit dans (35), définit bien une structure de graphe satisfaisant à $A1)$ $A2)$, i.e. permettant de définir une structure 2-polyédrale (qui sera automatiquement régulière \deleted{\unsure{)}}, on trouve en $G_0^+$ \deleted{\unsure{sur|formant}} un groupe d'automorphismes transitif sur l'ens des arêtes orientées du dit polyèdre\ldots).

Considérons d'abord $A1$, soit donc \deleted{$a\in\unsure{\mathbb P_2|\mathcal P_2|\mathfrak P_2}\unsure{(|(S)}$} $a\in A_\lambda$, on peut écrire

\[
a=\{s,\lambda t\}\qquad\text{avec }\{s,t\}\in A_0
\]

Soit $v\in S$, donc $v$ de la forme $\lambda u$, on aura donc

\[
\begin{aligned}
\deleted{\{s,\lambda t,v\}}&\\[-2pt]
{\underset{\substack{\text{i.e. }\{s,v\}\in A_\lambda\\\{\lambda t,v\}\in A_\lambda}}{a\cup\{v\}\in F_\lambda}}
 &\Longleftrightarrow\quad
 \begin{gathered}\{s,\deleted{\unsure{v|u}}\lambda u\}\in A_{\deleted{\unsure{0}}\lambda}\\
 \{\lambda t,\lambda u\}\in A_\lambda\end{gathered}\\
 &\Longleftrightarrow\quad
 \begin{gathered}\{s,u\}\in A_0\\\{t,\lambda u\}\in A_0\end{gathered}
\end{aligned}
\]

donc la condition $A1$ pour $\deleted{\unsure{A}}\Gamma_\lambda$ s'écrit

\[
\begin{aligned}
&\deleted{\text{Pour }(s,t)\in\unsure{S^2|A_0^2},\ \text{tel que }\{s,t\}\in}\\
&\text{Pour deux sommets }s,t\text{ adjacents pour }P_0,\ \inserted{\text{et }\lambda\in Z',}\text{ il existe}\\
&\text{exactement deux sommets }u\text{ tels que }\deleted{\text{l'on ait}}\\
&u\text{ soit adjacent à }s,\quad\lambda u\text{ adjacent à }t,\ \text{i.e.}\\
&\operatorname{Card}\ \lambda.A_0(s)\cap A_0(t)=2\qquad \text{si }
 \begin{cases}\{s,t\}\in A_0\\\lambda\in Z'\end{cases}
\end{aligned}\tag*{$(40)$}
\]
% archive sheet number: 116
